\chapter{Implementace [draft]}
\label{chap:implementace}

Tato kapitola popisuje praktickou realizaci navrženého doporučovacího systému. Postupně představíme využité nástroje, datovou sadu a způsob jejího předzpracování. Následně rozebereme implementaci samotného offline modelu, logiku vysvětlovacího modulu a nakonec i podobu uživatelského rozhraní, které slouží pro interakci se systémem.

\section{Použité technologie a knihovny}
\label{sec:technologie}

Celý projekt je programován v jazyce Python. Pro práci s daty a maticovými operacemi jsme využili standardní knihovny Pandas a NumPy. Samotná architektura modelů kolaborativního filtrování (CFAE) a řídkých autoenkodérů (SAE) je postavena na frameworku PyTorch, který umožňuje efektivní trénování na grafických kartách (GPU).

Pro automatizované pojmenování neuronů a tvorbu textových vysvětlení bylo nutné zvolit vhodný velký jazykový model (LLM). Zatímco dřívější výzkumy pro tyto účely často spoléhaly na výpočetně nejtěžší modely, jakými byly ve své době například GPT-3.5-turbo či Claude 2 \cite{Zeng2025}, nedávné studie prokázaly, že pro úlohy sémantické interpretace a evaluace lze úspěšně využít odlehčené (tzv. \textit{lightweight}) modely \cite{ONeill2024, Thelwall2026}. V rámci inženýrského návrhu jsme proto zvažovali tři hlavní přístupy s ohledem na potřeby a limity akademického prostředí:

\begin{itemize}
    \item \textbf{Placené API služby (např. OpenAI GPT-4o-mini):} Tento přístup představuje současný průmyslový standard s vysokou spolehlivostí formátování (např. JSON). Cenově se jedná o velmi efektivní řešení s náklady 0,15 USD za milion vstupních a 0,60 USD za milion výstupních tokenů \cite{LLMStats2026}, pro rozsáhlé experimenty nicméně stále vyžaduje vyhrazený finanční rozpočet.
    \item \textbf{Lokální open-weights modely (např. rodina Gemma 3 nebo Llama 3):} Tyto modely lze stáhnout a provozovat zcela offline \cite{GooglePricing2026, Thelwall2026}. Jejich hlavní výhodou je nulová cena za API a absolutní ochrana soukromí uživatelských dat. Tento přístup však nebyl zvolen z důvodu vysokých nároků na paměť a výpočetní výkon lokálních grafických karet.
    \item \textbf{Bezplatné úrovně (Free Tier) komerčních API:} Vývojářské platformy jako Google AI Studio poskytují pro vybrané modely (např. rychlý model Gemini 2.5 Flash-Lite) velkorysý bezplatný přístupový tarif \cite{GooglePricing2026}. Daní za využití této úrovně je skutečnost, že poskytovatel může odeslaná data využívat pro zlepšování svých vlastních produktů, což je ovšem pro akademické experimenty s již veřejnými daty obvykle přijatelný kompromis.
\end{itemize}

Na základě tohoto srovnání a pilotního testování jsme se pro napojení naší implementace rozhodli využít právě bezplatný tarif platformy \textbf{Google AI Studio s modelem řady Gemini}. Toto řešení zcela eliminuje finanční náklady na zpracování a anotaci tisíců neuronů ze SAE. Zmíněný kompromis v podobě sběru dat ze strany poskytovatele je v kontextu této práce akceptovatelný, jelikož pracujeme s veřejně dostupným datasetem recenzí z platformy Yelp a nedochází tak k žádnému narušení soukromí reálných uživatelů.


Samotné uživatelské rozhraní a webová aplikace pro následnou uživatelskou studii [draft: pak běží nad frameworkem Streamlit. Tato knihovna v jazyce Python byla zvolena proto, že umožňuje rychlé a efektivní vytvoření interaktivních ovládacích prvků (jako jsou posuvníky pro mechanismus \textit{steeringu}) a jejich přímé napojení na backendové výpočty modelu. Výrazně tak urychluje vývoj iterativních prototypů bez nutnosti psát složitý front-endový kód v jazyce JavaScript.]

\section{Předzpracování a reprezentace dat}
\label{sec:data}

K trénování a testování systému jsme zvolili veřejně dostupný dataset Yelp. Obsahuje velké množství interakcí (recenzí) mezi uživateli a body zájmu (POI), a především bohatá metadata potřebná pro pozdější vysvětlování. 

Data bylo nutné nejprve vyčistit. Odstranili jsme neaktivní uživatele a lokace s malým počtem hodnocení pomocí tzv. $K$-core filtrování (pro $K = 5$). [PŘÍPADNĚ: V případech, kdy byla data pro určité scénáře příliš řídká, jsme provedli syntetická doplnění uživatelských profilů, abychom simulovali konzistentnější historii návštěv.]

Z původních dat jsme extrahovali a zpracovali následující atributy:
\begin{itemize}
    \item \textbf{Geolokace:} Souřadnice jednotlivých míst (zeměpisná šířka a délka).
    \item \textbf{Kategorizace a tagy:} Každé místo v datasetu Yelp má přiřazenou sadu kategorií (např. \textit{Cafe, Mexican, Gym}). Tyto kategorie jsme převedli do textové podoby a využili je jako základní slovník pro sémantické mapování neuronů.
    \item \textbf{Interakce:} Samotná hodnocení (hvězdičky) jsme pro účely základního modelu převedli na implicitní zpětnou vazbu, kde hodnota 1 značí, že uživatel dané místo navštívil a ohodnotil alespoň 3 hvězdičkami.
\end{itemize}

Po předzpracování obsahuje dataset [DOPLŇ ČÍSLO] uživatelů, [DOPLŇ ČÍSLO] bodů zájmu a [DOPLŇ ČÍSLO] interakcí.

\section{Implementace modelu}
\label{sec:model}

Základní doporučovací algoritmus se skládá ze dvou částí, které trénujeme postupně. Jako první jsme naimplementovali základní model kolaborativního filtrování. [ZDE KRÁTCE POPIŠ, JESTLI POUŽÍVÁŠ ELSA NEBO NĚCO JINÉHO]. Tento model analyzuje implicitní interakce a každému uživateli přiřadí hustý (dense) vektor v latentním prostoru. Model byl optimalizován pomocí [DOPLŇ LOSS FUNKCI, např. binární křížové entropie].

Na výstup tohoto modelu jsme následně napojili vrstvu řídkého autoenkodéru (SAE). Abychom dosáhli požadované rozpletenosti (disentanglement) a vynutili aktivaci jen několika málo neuronů současně, použili jsme aktivační funkci Top-$K$. Ta při dopředném průchodu ponechá nenulových pouze $K$ neuronů s nejvyšší hodnotou a zbytek vynuluje. Během experimentů jsme testovali různé konfigurace a velikosti latentního prostoru. Ve finální verzi model pracuje s [DOPLŇ VELIKOST, např. 2048] dimenzemi a hodnotou $K$ nastavenou na [DOPLŇ ČÍSLO].

[DOPLŇ VIZUALIZACI]

\section{Systém vysvětlování a úprav}
\label{sec:explain_steer}

Aby byl model srozumitelný, museli jsme vyřešit překlad mezi natrénovanými číselnými hodnotami v SAE a textem. O to se stará vysvětlovací a řídící (Explain/Steer) modul.

Pro mapování neuronů na lidsky pochopitelné pojmy jsme využili jazykový model (LLM). Systém pro každý aktivní neuron vyhledá lokace, které ho nejvíce stimulují, a jejich metadata (kategorie a tagy) pošle jako prompt do LLM. Jazykový model z těchto dat vyvodí společný jmenovatel a vygeneruje pro daný neuron krátký popisný štítek (např. \textit{„Klidná místa pro práci“}). Stejný modul se pak podílí i na generování celkových textových vysvětlení k výsledným doporučením, kdy spojuje aktivní štítky uživatele s vlastnostmi doporučeného místa.

Součástí tohoto modulu je také backendová logika pro interakci (steering). Pokud uživatel vyjádří požadavek na změnu preferencí, modul vezme jeho aktuální řídký vektor a provede interpolaci s vektorem upravovaného neuronu. Takto modifikovaný profil následně pošle zpět do dekodéru, který bez nutnosti přetrénování bleskově vrátí nová doporučovací skóre.

\section{Uživatelské rozhraní}
\label{sec:ui}

Celá výše popsaná logika je zabalena do webové aplikace, se kterou pracují účastníci uživatelské studie. Rozhraní záměrně skrývá veškerou matematickou složitost modelu a prezentuje data vizuálně.

V levém panelu uživatel vidí svůj vygenerovaný profil. Ten je reprezentován sadou posuvníků, které odpovídají aktivním neuronům s přiřazenými štítky z LLM. Uživateli se po načtení stránky zobrazí úvodní seznam doporučených bodů zájmu. U každého doporučení je kromě názvu a kategorie uvedeno také krátké vysvětlení, proč systém dané místo vybral.

Hlavním prvkem rozhraní je možnost s posuvníky volně manipulovat. Jakýkoliv posun vyvolá požadavek na backend, systém přepočítá doporučení podle nové váhy faktorů a aktualizuje seznam míst na obrazovce. Tento okamžitý vizuální dopad (feedback loop) uživateli jasně ukazuje, že má nad algoritmem kontrolu.